% ============================================================
% 术语对照表（Glossary）模板
% 中文翻译的唯一术语标准
% 全文必须统一使用，不得出现同义异译
% 
% 使用方式：复制此文件到目标 paper-cn/ 目录，
% 根据具体论文修改术语条目。
% ============================================================

% --- 核心数学对象 ---
% half-derivation                 1/2-导子
% N_n                             幂零李代数 N_n（严格上三角矩阵）
% HalfDer(N_n)                    1/2-导子空间
% Hom(N_n, N_n)                   线性映射空间
% strictly upper-triangular       严格上三角
% nilpotent Lie algebra           幂零李代数
% commutator bracket              交换子括号
% Kronecker delta                 克罗内克 δ
% standard basis                  标准基
% maximal abelian ideal           极大交换理想
% I (ideal)                       I（理想）
% dim = n+5                       维数 = n+5
% dim = n error                   dim=n 错误

% --- 算子与空间 ---
% Psi operator                    Ψ 算子（half-Leibniz 条件）
% half-Leibniz morphism           半-Leibniz 态射
% half-Leibniz condition          半-Leibniz 条件
% bracket                         括号 [·,·]
% bracket-source term             括号-源项（bracketSource, BS）
% bracket-left term               括号-左项（bracketLeft, BL）
% bracket-right term              括号-右项（bracketRight, BR）
% bracket identity                括号恒等式（bracketIdentity, BI）
% adjoint representation          伴随表示 ad
% evaluation map                  求值映射 ev
% currying                        Currying（参数化）
% constraint operator             约束算子
% evaluation operator             求值算子
% vector equation space           向量方程空间 VEq
% scalar equation space           标量方程空间 SEq
% coordinate projection           坐标投影
% scalar linear functional        标量线性泛函
% image of ...                    ...的像
% kernel of ...                   ...的核
% rank-nullity theorem            秩–零度定理
% commutative diagram             交换图

% --- 证明架构：约束归约层 ---
% constraint reduction            约束归约
% chain bridge                    链桥（L2）
% width reduction                 宽度归约（L3）
% diagonal propagation            对角传播（L4）
% image restriction               像限制（L5）
% endpoint coupling               端点耦合（L6）
% cross-endpoint constraint       跨端点约束
% parameter count                 参数计数（L7）
% channel decomposition           信道分解

% --- 生成元分类 ---
% generator / generating family   生成元 / 生成族
% structural classification       结构分类
% T0 family                       T0 族（自配对）
% T1 family                       T1 族
% T2 family                       T2 族（T1 的 σ-镜像）
% T3 family                       T3 族（交换对）
% non-commuting pair              非交换对
% commuting pair                  交换对
% canonical decomposition         标准分解

% --- 导子分类 ---
% identity derivation             恒等导子
% simple-root center map          单根中心映射
% boundary cocycle                边界 cocycle
% centered derivation             中心化导子（对角为零）
% diagonal entry                  对角分量
% off-diagonal component          非对角分量
% identity component              恒等分量
% centered part                   中心化部分

% --- 信道变量 ---
% a-channel                       a-信道
% b-channel                       b-信道
% c-channel                       c-信道
% channel variable                信道变量
% channel restriction             信道限制
% self-pair constraint            自配对约束

% --- 结构概念 ---
% Dynkin diagram                  Dynkin 图
% Dynkin endpoint                 Dynkin 端点
% Dynkin involution               Dynkin 对合 σ
% sigma-equivariance              σ-等变性
% simple root                     单根
% non-local commuting pair        非局部交换对

% --- 矩阵相关 ---
% evaluation matrix               评价矩阵
% row / column                    行 / 列
% rank                            秩

% --- 上下界 ---
% lower bound                     下界
% upper bound                     上界
% linearly independent            线性无关
% explicit basis                  显式基

% --- 坐标约定 ---
% coordinate label                坐标标签
% coordinate convention           坐标约定
% scalar parameter                标量参数
% free parameter                  自由参数

% --- Lean 对应术语（保留原名，不翻译）---
% Lean formalization              Lean 形式化
% reference implementation        参考实现
% formal verification             形式化验证
% correspondence principle        对应原则

% --- 不翻译的代码标识符 ---
% coeffOf_cond                    保留
% coeffOf_map                     保留
% finrank_halfDer                 保留
% all_diag_equal                  保留
% boundary_rigidity               保留
% offdiag_zero_v_lt_n             保留
% bracketSource                   保留
% bracketIdentity                 保留
% bracketLeft                     保留
% bracketRight                    保留
% half_leibniz                    保留
% IsValidSupp                     保留
% I_filtered                      保留
% kerPhi                          保留

% --- 通用数学术语 ---
% characteristic                  特征
% field                           域
% linear map                      线性映射
% bilinear map                    双线性映射
% trilinear map                   三线性映射
% linear extension                线性扩张
% vanish / vanishes               为零 / 消失
% non-zero                        非零
% identically zero                恒为零
% invertible                      可逆
% uniquely determined             唯一确定
% support                         支撑集
% span                            张成空间
% basis element                   基元素
% decomposition                   分解
% induction                       归纳
% induction hypothesis            归纳假设
% base case                       基情形
% inductive step                  归纳步骤
% verification                    验证
% direct computation              直接计算
% expansion                       展开
% cancels / cancellation          抵消
% index mismatch                  指标不匹配
% structural omission             结构性遗漏
% structural property             结构性质
% consequence                     推论
% the exceptional case            例外情形
% counterexample                  反例

% --- 论文结构 ---
% Section                         节
% Subsection                      小节
% Appendix                        附录
% Organization of the paper       论文组织
% Motivation                      动机
% distinguished basis             特异基
% forward reference                前向引用
% proof                           证明
% Remark                          注记
% Corollary                       推论
% Definition                      定义
